\section{Quadratic characters with arbitrary translating centres}
\label{sec:quadratic}

We retain the sets $A,D=-B$, the partitions $S_p,S_p^c$, and the
densities $\sigma_p=|S_p|/p$ from Section~\ref{sec:preliminaries}.
For an odd prime $p$, let $\chi_p$ be the quadratic character of
$\F_p$, extended by zero at zero.  The translating residues in the
following proposition need not arise from a common integer or rational
number.

\begin{proposition}\label{prop:quadratic}
Under the assumed eventual decomposition of the primes,
\begin{equation}\label{eq:src3}
 \sum_{T\leq \log p\leq 2T}\frac{\log p}{p}
 \max_{t\in\F_p}
 \left|\E_{x\in S_p}\chi_p(x-t)\right|=o(T)
 \qquad (T\longrightarrow\infty).
\end{equation}
\end{proposition}

The main difficulty is that the translating residues are initially
unrelated.  An amplified moment and Poisson summation, with estimates
uniform over a polynomial-sized family of coefficient arrays, first
produce a common rational centre for many primes.  The quadratic large
sieve then confines positive fractions of the two endpoint tails to a
few quadratic kernels, whose populations contradict the collision estimate.

\pdfbookmark[2]{Biased block and amplified moment}{quadratic.moment}
\begin{proof}
\textbf{A biased prime block and the choice of scales.}
Suppose that \eqref{eq:src3} fails.  Mertens' estimate gives
\[
 \sum_{T\leq\log p\leq2T}\frac{\log p}{p}=T+O(1).
\]
Since each bias is at most one, there are a constant $\delta>0$ and an unbounded sequence of
$T$ for which the primes with maximal bias at least $\delta$ have
$\log p/p$-weight at least $c_\delta T$.
Throughout this proof all limiting assertions refer to this sequence,
and all constants may depend on this fixed bias.  Set
\[
 \rho=\frac35,\qquad \mu_0=10^{-6},\qquad \gamma=10^{-7},\qquad
 X=\exp(T^{1+\rho}),
\]
and define
\[
 A_X=A\cap[X^{9/10},X],\qquad
 D_X=D\cap[-X,-X^{9/10}].
\]
By Lemma~\ref{lem:size}, both tails have size
$\gg \sqrt X/(\log X)^3$.  Their uniform measures satisfy
Lemma~\ref{lem:collision}, for example with its fixed exponent $b=4$.
The prime bands used below lie below the corresponding cutoff
$\sqrt X/(\log X)^5$.

The nonnegative density term in \eqref{eq:src2} shows that primes with
$\sigma_p\notin[1/3,2/3]$ have total $\log p/p$-weight
$O(\log\log X)=O(\log T)$ in either of the bands considered here.
For any functions $\phi_p:\F_p\to\C$ with $|\phi_p|\leq1$, the same equation and
Cauchy--Schwarz give
\[
 \sum_{p\leq \sqrt X/(\log X)^5}\frac{\log p}{p}
 \left|\E_{x\in A_X}\phi_p(x)
       -\E_{x\in S_p}\phi_p(x)\right|^2
 \ll\log\log X,
\]
and the analogous assertion for $D_X$ and $S_p^c$.  This estimate
allows the test function to be chosen separately at every prime.
For each biased prime choose a maximizing translate and an orientation
$\epsilon_p\in\{1,-1\}$.  After deleting weight $O_\delta(\log T)$,
both empirical means are within $\delta/8$ of the corresponding
uniform means, and $1/3\leq\sigma_p\leq2/3$.
The two uniform means are related by
\[
 \E_{S_p^c}\chi_p(x-t)
   =-\frac{\sigma_p}{1-\sigma_p}\E_{S_p}\chi_p(x-t),
\]
because the complete character sum is zero.
Subdivision into intervals $[Z,2Z]$ therefore yields a set
$\mathcal P\subset[Z,2Z]$ of primes, with
$\log Z\asymp T$ and $J=|\mathcal P|\gg_\delta Z/\log Z$, such that
\begin{equation}\label{eq:src4}
 \E_{x\in A_X}\epsilon_p\chi_p(x-t_p)\geq c_\delta,
 \qquad
 \E_{x\in D_X}\epsilon_p\chi_p(x-t_p)\leq-c_\delta
 \quad(p\in\mathcal P),
\end{equation}
where $c_\delta>0$ is fixed.  Indeed, the original band contains
$O(T)$ such intervals, and one retains a positive constant of weighted
mass in at least one interval; on it each prime has weight
$O(\log Z/Z)$.

Apply the density part of \eqref{eq:src2} also to
$T^\gamma\leq\log p\leq2T^\gamma$.  Its full weight is
$T^\gamma+O(1)$, while the exceptional weight is $O(\log T)$.
Consequently some block $[z,2z]$, with $\log z\asymp T^\gamma$,
contains $\gg z/\log z$ primes for which
$1/3\leq\sigma_p\leq2/3$.  Put
\[
 K=\left\lfloor .02\frac{\log Z}{\log(2z)}\right\rfloor,
\]
and let $L$ be a product of $K$ distinct such primes.  There are enough
primes to do this, since $z/\log z$ exceeds every fixed power of $T$.
Let $k$ be the smallest even integer at least
$\log X/\log Z+10$.  Thus
\begin{equation}\label{eq:quad-scales}
 K\asymp T^{1-\gamma},\quad L\leq Z^{.02},\quad
 k\asymp T^\rho,\quad K=o(\log Z),\quad
 \frac Kk\gg T^{1-\gamma-\rho}.
\end{equation}
In particular all factors of $L$ are distinct from the primes in
$\mathcal P$.

\paragraph{An amplified moment and removal of repeated primes.}
For $\ell\mid L$ prime write
$f_\ell=\ind_{S_\ell}-\sigma_\ell$, periodically on $\Z$, and set
\[
 \lambda=\frac1{16},\qquad
 W(n)=\prod_{\ell\mid L}(1+\lambda f_\ell(n)),\qquad
 H(n)=\frac1J\sum_{p\in\mathcal P}\epsilon_p\chi_p(n-t_p).
\]
Fix an even nonnegative Schwartz function $\psi$ that is bounded below
by a positive constant on $[0,1]$ and has smooth compactly supported
Fourier transform, with convention
$\widehat\psi(y)=\int_\R\psi(x)e(-xy)\,dx$.
Such a function is obtained by squaring the inverse Fourier transform
of a sufficiently narrow real even smooth bump.
The weight $W$ is positive, has period $L$ and mean one, and satisfies
$W(n)\geq(1+\lambda/3)^K$ on $A_X$.
Since $k$ is even, Jensen's inequality and \eqref{eq:src4} give
\begin{equation}\label{eq:quad-moment-lower}
 I:=\sum_{n\in\Z}\psi(n/X)W(n)H(n)^k
   \gg \frac{\sqrt X}{(\log X)^3}(1+\lambda/3)^Kc_\delta^k
   \geq\sqrt X\exp(.016K)
\end{equation}
for sufficiently large $T$.  Here
$\log(1+\lambda/3)>.0206$, and $k+\log\log X=o(K)$.
Poisson summation, in its periodized form
\cite[Appendix~D.2]{MontgomeryVaughan2007}, gives for the $L$-periodic
function $W$
\begin{equation}\label{eq:quad-total-weight}
 I_0:=\sum_n\psi(n/X)W(n)=X\widehat\psi(0)\ll X:
\end{equation}
all nonzero frequencies vanish once $X/L$ exceeds the fixed support
radius of $\widehat\psi$.

We shall use the following elementary bound.  If
$H=J^{-1}\sum_{i=1}^J y_i$ with $y_i\in\{0,1,-1\}$ and $k$ is even,
then, for an absolute constant $C$,
\begin{equation}\label{eq:src5}
 \left|H^k-\frac1{J^k}
       \sum_{\substack{i_1,\ldots,i_k\in[J]\\\text{distinct}}}
                   y_{i_1}\cdots y_{i_k}\right|
 \leq\sum_{1\leq j\leq k/2}
             \left(\frac{k^C}{J}\right)^j|H|^{k-2j}.
\end{equation}
To prove it, express the distinct-index sum by M\"obius inversion on
set partitions; see \cite[Section~7, Example~1]{Rota1964}.
Equivalently, sum over permutations of $[k]$, with a
cycle of length $m$ contributing $\sum_i y_i^m$ and with the usual
permutation sign.  A partition block of size $m$ has coefficient
$(-1)^{m-1}(m-1)!$, since precisely $(m-1)!$ cycles have that support.
If a permutation has $k-2j$ odd cycles, those cycles contribute
$(JH)^{k-2j}$, whereas every even cycle contributes at most $J$ in
absolute value.  The total number of cycles is at most $k-j$.
Moreover, at most $3j$ positions belong to cycles of length greater
than one: an odd cycle of length $m\geq3$ uses $m$ positions while
contributing $m-1$ to $2j$, and an even cycle contributes all its
positions to $2j$.  There are at most $k^{O(j)}$ permutations with
this many moved positions.  The identity permutation is the term
$H^k$; all the other terms give \eqref{eq:src5}.

Apply \eqref{eq:src5} pointwise and then use H\"older with respect to
the positive measure $\psi(n/X)W(n)$.  The relative error in $I$ is
at most
\[
 \sum_{j=1}^{k/2}
 \left(\frac{k^C}{J}\left(\frac{I_0}{I}\right)^{2/k}\right)^j=o(1).
\]
Indeed, the ratio is at most a constant times
$k^C\log Z\exp(-.032K/k)$, because $X^{1/k}\leq Z$;
\eqref{eq:quad-scales} makes this tend to zero.

\pdfbookmark[2]{Poisson summation and uniform moments}{quadratic.poisson}
\paragraph{Poisson summation with arbitrary translations.}
Henceforth choose an ordered tuple of $k$ distinct primes of
$\mathcal P$ uniformly at random.  Its product is denoted by $M$.
Expectation with respect to these tuples will also be written
$\E_M$, since the functions in use depend only on their product.
Let
\[
 \chi_M=\prod_{p\mid M}\chi_p,\qquad
 0\leq t_M<M,\quad t_M\equiv t_p\pmod p\quad(p\mid M),
\]
and set
\[
 R=M/X,\qquad \theta=t_M/M,\qquad
 h_0\equiv-t_M\overline M\pmod L,\quad 0\leq h_0<L.
\]
The bars here and below indicate inverses modulo the modulus in the
expression.  The choice of $k$ gives, uniformly in the tuple,
\begin{equation}\label{eq:quad-R-range}
 Z^{10}\leq R\leq Z^{12+o(1)}.
\end{equation}
For $d\mid L$ put
\[
 f_d(x)=\prod_{\ell\mid d}f_\ell(x),\qquad f_1=1,\qquad
 g_d(b)=\sqrt d\,\E_{x\bmod d} f_d(x)e_d(-bx).
\]
The product is interpreted by the Chinese remainder theorem.  Fourier
orthogonality and the same theorem imply
\begin{equation}\label{eq:quad-g-properties}
 \E_{b\bmod d}|g_d(b)|^2\leq1,\quad |g_d(b)|\leq\sqrt d,\quad
 g_d(b)=\prod_{\ell\mid d}
              g_\ell\bigl(b\overline{d/\ell}\bigr).
\end{equation}
In particular $g_d$ vanishes on nonunits when $d>1$, since each
$f_\ell$ has mean zero.  We use $g_1=1$.

Expand $W=\sum_{d\mid L}\lambda^{\omega(d)}f_d$.
For a fixed $d$, Poisson summation modulo $Md$ gives the complete
transform of $f_d(n)\chi_M(n-t_M)$.  Its factor modulo $d$ is
$\sqrt d\,g_d(-u\overline M)$, and its factor modulo $M$ is
\[
 \tau(\chi_M)\chi_M(u)\chi_M(d)
       e_M(t_Mu\overline d),
\]
where $|\tau(\chi_M)|=\sqrt M$ by the primitive quadratic Gauss-sum
identity \cite[Theorems~9.6--9.7]{MontgomeryVaughan2007}.
Since $t_M+Mh_0$ is divisible by every $d\mid L$,
\[
 e_M(t_Mu\overline d)=e\bigl((\theta+h_0)u/d\bigr).
\]
It follows that the transform of the distinct-prime contribution,
after division by $\sqrt X$ and multiplication by a unit complex
number depending only on the tuple, is
\begin{equation}\label{eq:src6}
 \sum_{d\mid L}\frac{\lambda^{\omega(d)}\chi_M(d)}{\sqrt{Rd}}
 \sum_{u\in\Z}\chi_M(u)g_d(-u\overline M)
       e\bigl((h_0+\theta)u/d\bigr)
       \widehat\psi\bigl(u/(Rd)\bigr).
\end{equation}
The orientations $\prod_{p\mid M}\epsilon_p$ are included in that
unit.  The term $u=0$ vanishes.  Since $f_d$ and $\psi$ are real,
the negative-frequency part is $\chi_M(-1)$ times the conjugate of
the positive-frequency part.  Write the latter as $\mathcal F_M$.
The normalization from sampling distinct tuples is
$(J)_k/J^k=1+o(1)$, where $(a)_r=a(a-1)\cdots(a-r+1)$.
The moment bound and the repeat estimate therefore imply
\begin{equation}\label{eq:quad-fourier-lower}
 \E_M|\mathcal F_M|\geq\exp(.015K).
\end{equation}
For example, the absolute value of \eqref{eq:src6} is at most
$2|\mathcal F_M|$, while the averaged original distinct contribution
is $(1-o(1))I/\sqrt X$.

Write uniquely
$u=svw^2>0$, where $s$ is squarefree and coprime to $L$, $v\mid L$,
and $w\geq1$.  For a positive integer $P$ define
\begin{equation}\label{eq:src7}
 \begin{aligned}
 N_{dsv}&=\sqrt{Rd/(sv)},\\
 G_d^{(P)}(s,v)&=\frac1{N_{dsv}}
   \sum_{\substack{w\geq1\\P\mid w}}
     g_d(-svw^2\overline M)
     e\bigl((h_0+\theta)svw^2/d\bigr)
     \widehat\psi\bigl(svw^2/(Rd)\bigr),\\
 C_s^{(P)}&=\sum_{v\mid L}\frac{\chi_M(v)}{\sqrt v}
            \sum_{d\mid L}\lambda^{\omega(d)}\chi_M(d)
                                G_d^{(P)}(s,v).
 \end{aligned}
\end{equation}
Since $\chi_M(w^2)=\ind_{(w,M)=1}$, inclusion-exclusion gives
\begin{equation}\label{eq:src8}
 \mathcal F_M=
  \sum_{P\mid M}(-1)^{\omega(P)}
    \sum_{\substack{s\leq Z^{14}\\s\text{ squarefree}\\(s,L)=1}}
       \frac{\chi_M(s)}{\sqrt s}\,C_s^{(P)}.
\end{equation}
The fixed support of $\widehat\psi$ and
\eqref{eq:quad-R-range} ensure $u\ll RL<Z^{14}$ and $w<Z^7$.
Consequently only $P\leq Z^7$ occur in this expression.  A divisor
of $M$ this small contains at most seven prime factors, so the number
of outer terms is $O(k^7)$.

To relate the translating residues, we shall force a large value of
$G_d^{(1)}(s,v)$ with $s<L^4$ for sufficiently many prime tuples.
On a progression $w=x+dj$, its remaining quadratic phase has coefficient
$svd\theta$, so a large value will give a small-denominator approximation
to $\theta=t_M/M$.  Since the arrays $C_s^{(P)}$ depend on the sampled
tuple $M$, the intervening moment estimate must be uniform over a fixed
family containing their discretizations.

\paragraph{A moment estimate for a polynomial-sized family of arrays.}
Put
$l=\lfloor T^{\mu_0}\rfloor$ and $u_0=2l$.
Let $\mathcal B$ be any fixed family of complex arrays $b=(b_s)$,
indexed by the squarefrees $s\leq Z^{14}$ coprime to $L$.
Suppose that, for a fixed constant $c$,
$|\mathcal B|\leq Z^c$ and $\sum_s|b_s|\leq Z^c$ for every
$b\in\mathcal B$.  We claim
\begin{equation}\label{eq:src9}
 \begin{aligned}
 &\left(\E_M\max_{b\in\mathcal B}
           \left|\sum_s\chi_M(s)b_s\right|^{2l}\right)^{1/(2l)}
 \\
 &\qquad\leq \exp(o(K))\max_{b\in\mathcal B}
           \left(\sum_su_0^{\omega(s)}|b_s|^2\right)^{1/2}
       +O(Z^{-10}).
 \end{aligned}
\end{equation}
The $o(K)$ is uniform for families with this fixed $c$.

To prove the claim, first replace the maximum of the moments by their
sum.  Each product $M$ has probability
$k!/(J)_k\leq2k!/J^k$, so positivity permits domination by the sum
over all odd $m\leq(2Z)^k$ with $\chi_M(s)$ replaced by $(s/m)$.
Expand one $2l$-th moment.  If the product $n$ of its $2l$ indices is
nonsquare, $m\mapsto(n/m)$ on odd $m$, extended by zero on even
$m$, is a nonprincipal quadratic Dirichlet character with possible
additional zero factors.  It has a period at most $8n$ and mean zero
over a period.  One can see this by using the nontrivial squarefree
kernel of $n$ and then imposing coprimality to its square part;
the induced character remains nonprincipal on the units of a modulus
dividing $8\operatorname{rad}(n)$.
As $n\leq Z^{28l}$, its sum over any initial interval is
$O(Z^{28l})$.  The total absolute nonsquare error for an array is
therefore at most $Z^{28l}(\sum_s|b_s|)^{2l}=Z^{O_c(l)}$.

When the product of the indices is square, its Jacobi symbol is either
zero or one.  Its contribution in absolute value is bounded by
$(2Z)^k$ times the $2l$-th moment of
\[
 F(\varepsilon)=\sum_s|b_s|\prod_{p\mid s}\varepsilon_p,
\]
where the $\varepsilon_p$ are independent uniform signs.
Indeed, the expectation of a product of these monomials is one
exactly when the product of the corresponding squarefree indices is
a square, and is zero otherwise.
We use the even-moment form of Bonami's hypercontractive inequality
\cite[Chapter~III, Theorems~2--3]{Bonami1970} and include its elementary
proof here.
For real $x,y$ the binomial theorem gives
\[
 \left(\E_{\varepsilon=\pm1}|x+\varepsilon y|^{2l}\right)^{1/l}
 \leq x^2+(2l-1)y^2,
\]
because
$\binom{2l}{2i}\leq\binom li(2l-1)^i$.
Apply this one sign at a time.  Minkowski's inequality, applied to
the square functions of the remaining signs, shows inductively that
\[
 \|F\|_{2l}^2\leq
       \sum_s(2l-1)^{\omega(s)}|b_s|^2
 \leq\sum_su_0^{\omega(s)}|b_s|^2.
\]
This is the tensorized sign-moment form of Bonami hypercontractivity;
see \cite[Chapter~III, Theorem~3, p.~376]{Bonami1970}. The preceding
argument proves the needed form directly; the arithmetic moment transfer
around it is the adaptation used here. Squarefreeness ensures that each
sign occurs with degree at most one in every monomial.

The $2l$-th root of the square-product transfer factor is at most
\[
 \left(Z^c\frac{2k!}{J^k}(2Z)^k\right)^{1/(2l)}
 \leq\exp\left(O_c\left(\frac{T+k\log T}{l}\right)\right)
 =\exp(o(K)).
\]
The last assertion uses $\mu_0>\gamma$ and $\rho<1$.
The nonsquare error, after the same probability factor and taking a
$2l$-th root, is at most
\[
 \exp\left(-\frac{k\log Z}{2l}
            +O_c(\log Z)+O\left(\frac{k\log T}{l}\right)\right),
\]
which is $O(Z^{-10})$ since $l=o(k)$.  This proves
\eqref{eq:src9}.

\paragraph{Discretizing the coefficients and removing $P\ne1$.}
The coefficients in \eqref{eq:src7}, except for their dependence on
$\theta,R,P$, are determined by $M\bmod4L$ and $h_0\bmod L$.
For the quadratic symbols of divisors of $L$, this follows from
quadratic reciprocity; the inverse of $M$ modulo every $d\mid L$
is already determined modulo $L$.
We allow every unit residue $M\bmod4L$, every $h_0\bmod L$,
every integer $1\leq P\leq Z^7$, and grids of mesh at most
$Z^{-100}$ for
\[
 0\leq\theta\leq1,\qquad Z^{10}\leq R\leq2Z^{13}.
\]
This is a family of $Z^{O(1)}$ parameter choices, with an absolute
constant exponent.  The larger range for $R$ also accommodates
nearest grid points.

These grids approximate the coefficients uniformly.  On the support
of a summand in $G_d^{(P)}$ one has $w\ll N_{dsv}$ and
$svw^2/d\ll R$.  The number of positive multiples of $P$ in this
range is $O(N_{dsv}/P)$, with no extra term required; if the range
contains any multiple at all, then $N_{dsv}/P$ is bounded below.
Differentiating the formula, and using the smooth compact support,
gives uniformly
\[
 |\partial_\theta G_d^{(P)}|\ll R\sqrt L/P,\qquad
 |\partial_R G_d^{(P)}|\ll \sqrt L/(RP).
\]
The derivative with respect to $R$ includes the factor $N_{dsv}^{-1}$
and the argument of $\widehat\psi$; both have the indicated bound.
Thus nearest grid replacement changes each $G_d^{(P)}(s,v)$ by
$O(Z^{-80})$.  The same estimates hold on the entire enlarged grid
range, where the support still has $s<Z^{14}$ and $w<Z^7$.
Even after summing the coefficient errors absolutely in
\eqref{eq:src8}, their contribution is $o(1)$: there are at most
$Z^{14}$ indices, divisor sums cost $\exp(O(K))=Z^{o(1)}$, and
there are $O(k^7)$ outer terms.
The crude bound $|G_d^{(P)}|\ll\sqrt L/P$ also shows that every
array $b_s=s^{-1/2}C_s^{(P)}$ has polynomial $\ell^1$ norm.

We shall repeatedly use
\begin{equation}\label{eq:src10}
 \begin{aligned}
 \sum_{v\mid L}v^{-1/2}&=1+o(1),\\
 \sum_{\substack{s\leq Z^{14}\\s\text{ squarefree}}}
       \frac{u_0^{\omega(s)}}s
 &\leq\prod_{p\leq Z^{14}}(1+u_0/p)\\
 &\leq\exp\left(u_0\sum_{p\leq Z^{14}}\frac1p\right)
 =\exp(o(K)).
 \end{aligned}
\end{equation}
For the first assertion, the logarithm of the divisor product is
$O(K/\sqrt z)=o(1)$.  For the last one it is
$O(T^{\mu_0}\log T)=o(K)$.
If $P\mid M$ and $P\ne1$, then $P\geq Z$.
For all arrays with this restriction, their weighted square norm in
\eqref{eq:src9} is at most
\[
 Z^{-1}\sqrt L\,(1+\lambda)^K\exp(o(K))
   =Z^{-1+.01+o(1)}.
\]
The maximum over all integer $P$ in the grid family bounds the
adaptive choice of the actual divisors.  Equation~\eqref{eq:src9}
and the $O(k^7)$ count show that all $P\ne1$ terms in
\eqref{eq:src8} have total expected absolute value $o(1)$.
We now set $P=1$ and suppress its superscript.

\pdfbookmark[2]{A small-kernel witness}{quadratic.witness}
\paragraph{A small-kernel quadratic sum must be large.}
We shall prove that with probability at least $\exp(-O(K))$ there
exist a squarefree $s<L^4$ coprime to $L$ and divisors $v,d\mid L$
such that
\begin{equation}\label{eq:src11}
 |G_d(s,v)|\geq2^{\omega(d)/2}e^{-K}.
\end{equation}
For $s<L^4$, on the entire grid range,
\[
 \frac{N_{dsv}}d=\sqrt{\frac R{svd}}\geq Z^{4.9}
\]
for sufficiently large $T$.
If $sv$ is nonunit modulo $d$, the function
$g_d(-svw^2\overline M)$ is identically zero.  Otherwise, because
$g_d$ vanishes at nonunits, the squaring map has fibers of size at
most $2^{\omega(d)}$ on every relevant residue.  Therefore
\[
 \E_{w\bmod d}|g_d(-svw^2\overline M)|^2\leq2^{\omega(d)},
 \qquad
 \E_{w\bmod d}|g_d(-svw^2\overline M)|\leq2^{\omega(d)/2}.
\]
The same assertions hold for $d=1$ by convention.
In an interval of length $O(N_{dsv})$, each residue occurs
$O(N_{dsv}/d)$ times.  Consequently
\begin{equation}\label{eq:quad-small-G-upper}
 |G_d(s,v)|\ll2^{\omega(d)/2}\qquad(s<L^4).
\end{equation}
By Minkowski and \eqref{eq:src10}, the corresponding small-$s$
arrays satisfy
\[
 \left(\sum_{s<L^4}\frac{u_0^{\omega(s)}}s|C_s|^2\right)^{1/2}
 \leq\exp(o(K))(1+\sqrt2\lambda)^K=\exp(O(K)).
\]
If \eqref{eq:src11} fails for an actual tuple, its nearest grid point
belongs, for sufficiently large $T$, to the fixed subfamily on which
every small-$s$ value is bounded by
$2^{\omega(d)/2}\exp(-.9K)$.  This follows since the mesh error
$O(Z^{-80})$ is negligible compared with $e^{-K}$.
The weighted square norm of each small-$s$ array in that subfamily is
at most
\begin{equation}\label{eq:quad-small-off-event}
 \exp\bigl((- .9+\log(1+\sqrt2\lambda)+o(1))K\bigr),
\end{equation}
which is exponentially small.

It remains to bound all large-$s$ arrays, uniformly in their
parameters, by $\exp((.006+o(1))K)$ in the weighted square norm.
Fix $v\mid L$ and a segment $S\leq s<2S$, where
$L^4\leq S\leq Z^{14}$, and put $\tau=.005$.
There are $O(\log Z)$ segments; their total cost in Minkowski's
inequality is $\exp(o(K))$, as is the sum over $v$ with the weights
$v^{-1/2}$ in \eqref{eq:src10}.

For the low-weight indices $u_0^{\omega(s)}\leq e^{\tau K}$,
bound the squared norm by $e^{\tau K}$ times the unweighted one.
By positivity this unweighted sum of squared moduli can be extended
to all integers in $[S,2S)$ before expanding its $d,d'$ cross terms;
the defining formula for $G_d(s,v)$ makes sense for these integers
as well.  We claim the correlation estimate
\begin{equation}\label{eq:src12}
 \left|\sum_{S\leq s<2S}
       \frac{G_d(s,v)\overline{G_{d'}(s,v)}}s\right|
 \ll\left(\frac{dd'}{(d,d')^2}\right)^{-1/2}.
\end{equation}
To verify it, expand the two sums over $w,w'$.  A nonzero summand
has
\[
 w\ll\sqrt{Rd/(Sv)},\qquad
 w'\ll\sqrt{Rd'/(Sv)}.
\]
Its periodic factor is
$F(s)=g_d(cs)\overline{g_{d'}(c's)}$, where
$c=-vw^2\overline M\pmod d$ and
$c'=-v(w')^2\overline M\pmod{d'}$.
If either scalar is a nonunit, the summand is identically zero.
Otherwise, on the period $q=[d,d']$, every normalized additive
Fourier coefficient of $F$ is bounded in modulus by
\[
 A_{d,d'}=\left(\frac{dd'}{(d,d')^2}\right)^{-1/2}.
\]
Indeed, at a prime in exactly one divisor, Fourier inversion of
$g_\ell$ gives an upper bound $1/\sqrt\ell$ since $|f_\ell|\leq1$.
At a common prime Cauchy--Schwarz and
\eqref{eq:quad-g-properties} give an upper bound one.
Chinese remaindering multiplies these local bounds.

The remaining phase is $e(\alpha s)$ for an arbitrary real number
$\alpha$, and the two transform factors form a smooth weight with
uniformly bounded supremum and total variation on $[S,2S]$.
For a $q$-periodic function whose Fourier coefficients are at most
$A_{d,d'}$, finite Fourier expansion and the geometric-sum bound give
\[
 \left|\sum_{s\in I}F(s)e(\alpha s)\right|
 \ll A_{d,d'}\bigl(|I|+q\log(2q)\bigr)
\]
for an interval $I$ of length at most $S$.  To see the uniformity in
$\alpha$, sum
$\min(S,\|\alpha+h/q\|^{-1})$ over $h\bmod q$: a closest
frequency costs at most $S$, and the others cost
$O(q\sum_{j\leq q}1/j)$.  Partial summation inserts the smooth
weight without changing this estimate.  Here $q\mid L$ and
$S\geq L^4$, so the bound is $O(A_{d,d'}S)$.
Finally, the factors $N_{dsv}^{-1}N_{d'sv}^{-1}/s$ equal
$v/(R\sqrt{dd'})$, independently of $s$, and there are at most
$O(R\sqrt{dd'}/(Sv))$ pairs $w,w'$.  This proves
\eqref{eq:src12}.
After the divisor sum, the low-weight squared norm is therefore
\begin{equation}\label{eq:quad-low-weight}
 \ll e^{\tau K}
       \prod_{\ell\mid L}(1+\lambda^2+2\lambda/\sqrt\ell).
\end{equation}

For completeness we give a residue-uniform estimate for the
high-weight indices.  For each residue class $a\bmod L$,
\begin{equation}\label{eq:quad-high-weight-mass}
 \sum_{\substack{S\leq s<2S, s\equiv a\ (L)\\
                  s\text{ squarefree}, (s,L)=1\\
                  u_0^{\omega(s)}>e^{\tau K}}}
                  u_0^{\omega(s)}
 \leq\frac S L e^{-10K}
\end{equation}
for sufficiently large $T$.  Given such an $s$, take $r$ to be the
product of its smallest $\lfloor\omega(s)/2\rfloor$ prime factors.
Then
\[
 r\leq\sqrt{2S},\quad(r,L)=1,\quad
 \omega(r)\geq\frac{\tau K}{2\log u_0}-1,\quad
 u_0^{\omega(s)}\leq u_0(u_0^2)^{\omega(r)}.
\]
There are at most $O(S/(Lr))$ multiples of $r$ in the specified class
and interval.  The usual additive constant in this count is absorbed
because $S/(Lr)\geq\sqrt S/(\sqrt2L)\gg1$.
Allowing all possible squarefree $r$, the desired sum is at most
$O(S/L)$ times
\[
 \begin{aligned}
 u_0\sum_{\substack{r\leq\sqrt{2S},\ r\text{ squarefree}\\
                 \omega(r)\geq\tau K/(2\log u_0)-1}}
       \frac{(u_0^2)^{\omega(r)}}r
 &\leq u_0T^{1/4-\tau K/(8\log u_0)}
        \prod_{p\leq\sqrt{2S}}(1+u_0^2T^{1/4}/p)\\
 &\leq u_0T^{1/4-\tau K/(8\log u_0)}
        \exp\bigl(O(u_0^2T^{1/4}\log T)\bigr)
 \leq e^{-20K}.
 \end{aligned}
\]
The last inequality follows since the negative main exponent is
$-(\tau/(8\mu_0)+o(1))K=-(625+o(1))K$, whereas the positive
exponent is $O(T^{1/4+2\mu_0}\log T)=o(K)$.
This proves \eqref{eq:quad-high-weight-mass}.

In the expansion used for \eqref{eq:src12}, sum the high-weight part
absolutely.  Its periodic factor satisfies
$\E_{s\bmod L}|F(s)|\leq1$, by Cauchy--Schwarz and
\eqref{eq:quad-g-properties}.  Thus
\eqref{eq:quad-high-weight-mass}, followed by the same normalization
and pair count, bounds its contribution by $O(e^{-10K})$ for each
pair of divisors.  The divisor weights sum to
$(1+\lambda)^{2K}$, which leaves an exponentially negligible bound.
In \eqref{eq:quad-low-weight} the logarithm of the product is
$K\log(1+\lambda^2)+o(K)$, because every $\ell\geq z$.
Combining the two parts, and then summing the segments and $v$, gives
\begin{equation}\label{eq:quad-large-s-norm}
 \left(\sum_{\substack{s\geq L^4\\s\text{ squarefree}\\(s,L)=1}}
             \frac{u_0^{\omega(s)}}s|C_s|^2\right)^{1/2}
 \leq\exp((.006+o(1))K),
\end{equation}
since $(\tau+\log(1+\lambda^2))/2<.006$.

Let $\mathcal E$ be the event in \eqref{eq:src11}, and split the
$P=1$ inner expression in \eqref{eq:src8} into small and large $s$.
Apply \eqref{eq:src9} to the grid families for the large part and to
the fixed subfamily in \eqref{eq:quad-small-off-event}.
Grid replacement costs $o(1)$ in each application.  It follows that
the large part has expected absolute value at most
$\exp((.006+o(1))K)$, and the small part on $\mathcal E^c$ has
exponentially small expectation.  In the latter assertion it is
important that the maximum over the whole fixed subfamily bounds
every tuple in $\mathcal E^c$; no conditional character estimate is
being asserted.
The whole small part has $L^{2l}$ probability norm $\exp(O(K))$ by
\eqref{eq:quad-small-G-upper} and \eqref{eq:src9}.
In view of \eqref{eq:quad-fourier-lower} and the negligible $P\ne1$
terms, its expected absolute value on $\mathcal E$ is at least
$\tfrac12\exp(.015K)$ for large $T$.
H\"older's inequality now implies
\begin{equation}\label{eq:quad-event-probability}
 \Prob(\mathcal E)\geq\exp(-O(K)),
\end{equation}
as claimed.

\pdfbookmark[2]{From a rational approximation to a common centre}{quadratic.centre}
\paragraph{Recovering a rational approximation from the witness.}
Fix an actual tuple in $\mathcal E$ and a witness $s,v,d$.
Split the $w$-sum in \eqref{eq:src7} into progressions modulo $d$.
The $g_d$ factor and the $h_0$ phase are constant on each
progression, and the sum of their absolute amplitudes over the
classes is at most $d\,2^{\omega(d)/2}$.
Put $Y_0=N_{dsv}/d$.  Equation~\eqref{eq:src11} shows that on some
progression $w=x+dj$ the remaining smooth quadratic sum has modulus
at least $Y_0e^{-K}$.  Its quadratic coefficient in $j$ is
$\alpha=svd\theta$.  The weight
$\widehat\psi(((j+x/d)/Y_0)^2)$ is supported on an interval of
length $O(Y_0)$ and has bounded supremum and total variation,
uniformly in the progression.  Partial summation therefore gives an
unweighted interval sum
\[
 \left|\sum_{j\in I}e(\alpha j^2+\beta j)\right|
       \gg Y_0e^{-K},\qquad |I|\ll Y_0,
\]
with an arbitrary real linear coefficient $\beta$.
We have $Y_0\gg Z^4$, $\log Y_0\asymp T$, and
$\log T=o(K)=o(T)$.

We record the inverse conclusion, including its dependence on the
length:
\begin{equation}\label{eq:src13}
 1\leq q\leq\exp(O(K)),\qquad
 \|q\alpha\|_{\R/\Z}\leq\exp(O(K))Y_0^{-2}
\end{equation}
for some integer $q$.
For a proof, let $Q_1=\lfloor Y_0^2e^{-8K}\rfloor$ and apply
Dirichlet approximation to $2\alpha$.  After reducing the resulting
fraction, we have coprime integers $b,r$, with
$1\leq r\leq Q_1$ and $|2r\alpha-b|\leq1/Q_1$.
In particular $|2\alpha-b/r|\leq1/r^2$.
One differencing step, in which the linear coefficient and the
interval location do not affect the absolute bound, gives
\[
 \left|\sum_{j\in I}e(\alpha j^2+\beta j)\right|^2
 \ll Y_0+\sum_{1\leq h\ll Y_0}
                 \min(Y_0,\|2h\alpha\|^{-1})
 \ll (Y_0/r+1)(Y_0+r\log(2r)).
\]
For the second bound split the shifts into blocks of length at most
$r/2$.  Within such a block two values of $2h\alpha$ are separated
modulo one by at least $1/(2r)$, because the reduced rational values
are separated by $1/r$ and the approximation error between them is
at most $1/(2r)$.  The sum of the displayed minima in one block is
$O(Y_0+r\log(2r))$.  Bounded $r$ is covered by the same estimate.
If $e^{8K}<r\leq Q_1$, the last bound is at most
\[
 O\bigl(Y_0^2e^{-8K}\log(2Y_0)+Y_0\log(2Y_0)\bigr)
       =o(Y_0^2e^{-2K}),
\]
contrary to the lower bound for the sum.
Thus $r\leq e^{8K}$, and $q=2r$ proves
\eqref{eq:src13}, since $1/Q_1\ll e^{8K}Y_0^{-2}$.

Let $a=qsvd$.  Then
\[
 1\leq a\leq L^6e^{O(K)}\leq Z^{.12+o(1)}.
\]
Choose an integer $b'$ nearest to $a\theta$ and put
$n=at_M-b'M$.  As $Y_0^{-2}=svd/R$, Equation~\eqref{eq:src13}
gives
\begin{equation}\label{eq:quad-integer-lift}
 |n|\leq aXe^{O(K)}\leq XZ^{.12+o(1)},\qquad
 n\equiv at_p\pmod p\quad(p\mid M).
\end{equation}
There are at most $Z^{.12+o(1)}$ choices of $a$.
Using \eqref{eq:quad-event-probability}, $(J)_k=(1-o(1))J^k$, and
$e^{O(K)}=Z^{o(1)}$, pigeonholing gives a fixed positive integer
$a\leq Z^{.12+o(1)}$ for which at least $J^kZ^{-.13}$ ordered
distinct tuples have a lift satisfying \eqref{eq:quad-integer-lift}.

\paragraph{One common centre for many primes.}
Put $H_0=XZ^{.15}$, so that $H_0\leq Z^{k-10+.15}$, and for every
integer $|n|\leq H_0$ let
\[
 r(n)=|\{p\in\mathcal P:n\equiv at_p\pmod p\}|.
\]
Counting ordered tuples and then counting lifts modulo products of
$k-10$ distinct primes gives
\begin{equation}\label{eq:quad-lift-factorials}
 \sum_{|n|\leq H_0}(r(n))_k\geq J^kZ^{-.13},\qquad
 \sum_{|n|\leq H_0}(r(n))_{k-10}\ll Z^{.15}J^{k-10}.
\end{equation}
For the second estimate each product is at least $Z^{k-10}$, so
the number of its lifts in the interval is $O(Z^{.15})$.
The first estimate counts at least one lift for every tuple already
obtained.  The terms with $r(n)<Z^{.6}$ contribute at most
\[
 Z^6\sum_n(r(n))_{k-10}\ll Z^{6.15}J^{k-10}
       =o(J^kZ^{-.13})
\]
to the first sum in \eqref{eq:quad-lift-factorials}.

For each remaining integer take its set of matching primes.  Two
different integers have at most $k$ common matching primes: a product
of $k+1$ such primes would divide their nonzero difference, whose
absolute value is at most $2H_0<Z^{k+1}$.
Let there be $b$ high sets, of total cardinality $R_0$.
If a prime occurs in $d_p$ of them, Cauchy--Schwarz and the
intersection bound give
\[
 \frac{R_0^2}{J}\leq\sum_pd_p^2\leq R_0+kb^2.
\]
Since $b\leq R_0/Z^{.6}$ and $kJ/Z^{1.2}=o(1)$, this implies
$R_0\leq2J$ for sufficiently large $T$.
If $J_0$ is the maximum size of a high set, then
\[
 \tfrac12J^kZ^{-.13}
 \leq\sum_{n:\,r(n)\geq Z^{.6}}(r(n))_k
 \leq J_0^{k-1}R_0\leq2JJ_0^{k-1}.
\]
Thus one integer $n$ has a matching set $\mathcal P_0$ satisfying
\begin{equation}\label{eq:quad-common-centre-size}
 J_0:=|\mathcal P_0|\geq JZ^{-O(1/k)}.
\end{equation}
Reduce $n/a=h/m$ to lowest terms, with $m>0$.
Since $a<Z^{.13}<p$ for $p\in\mathcal P_0$, reduction of this
identity modulo $p$ is valid and yields
\begin{equation}\label{eq:quad-common-centre}
 t_p\equiv h/m\pmod p\quad(p\in\mathcal P_0),\qquad
 m\leq Z^{.13}=X^{o(1)},\qquad |h|\leq XZ^{.15}.
\end{equation}
This is the first point at which the initially arbitrary centres
have been related to one another.

\pdfbookmark[2]{Quadratic kernels and disjoint supports}{quadratic.kernels}
\paragraph{The quadratic large sieve and the number of kernels.}
Prime-product character tests for squarefree kernels also appear in
Green and Harper~\cite[Lemmas~6.1--6.2]{GreenHarper}.
Here the preceding argument has first related the independently chosen
translating centers.
We use Heath-Brown's quadratic large sieve
\cite[Theorem~1]{HeathBrown1995}: for arbitrary complex coefficients
on odd positive squarefree $s\leq S$,
\begin{equation}\label{eq:quad-heath-brown}
 \sum_{u\leq U}^{*}
     \left|\sum_s b_s(s/u)\right|^2
 \ll_\varepsilon (US)^\varepsilon(U+S)\sum_s|b_s|^2,
\end{equation}
where the starred sum is over odd positive squarefree $u$.
The same bound, up to an absolute factor, holds for $(u/s)$ with
nonzero signed squarefree $u$ satisfying $|u|\leq U$.
To check this extension explicitly, write $u=\pm2^e v$, with
$e\in\{0,1\}$ and $v$ positive odd squarefree, and separate
$v\bmod4$.  Quadratic reciprocity turns $(v/s)$ into $(s/v)$ times
a sign depending only on $s$ and the fixed class of $v$.
The additional factors $(\pm2^e/s)$ can also be absorbed into
$b_s$.  The coefficient square norm is unchanged.  Dropping the
restriction on $v\bmod4$ by positivity and applying
\eqref{eq:quad-heath-brown} in each of these eight cases proves the
extension, including $u=\pm1$.

Put $\epsilon'_p=\epsilon_p\chi_p(m)$ for
$p\in\mathcal P_0$.  Equations~\eqref{eq:src4} and
\eqref{eq:quad-common-centre} show that
\[
 F(x)=\frac1{J_0}\sum_{p\in\mathcal P_0}
                          \epsilon'_p\chi_p(mx-h)
\]
has mean at least $c_\delta$ on $A_X$ and at most $-c_\delta$ on
$D_X$.  Since $|F|\leq1$, on a fixed positive fraction of each
tail we have $|F(x)|\geq c_\delta/2$.
On these subsets $mx-h\ne0$, and write uniquely
\[
 mx-h=u_xt_x^2,\qquad u_x\text{ signed squarefree},\quad t_x\geq1.
\]
Uniformly $|u_x|t_x^2\leq mX+|h|=X^{1+o(1)}$.
At most $O(k)$ primes of $\mathcal P_0$ can divide $t_x$, by taking
logarithms of their product.  Replacing $\chi_p(mx-h)$ by
$\chi_p(u_x)$ changes the average by $O(k/J_0)=o(1)$.
Thus the average with $u_x$ remains bounded away from zero.

Raise that average to the even power $k$ and apply
\eqref{eq:src5}, now with $J_0$ in place of $J$.
Its relative error is $o(1)$, since the mean has modulus at least a
fixed positive constant and $k^C/J_0=o(1)$.
We obtain a single coefficient array, the same for every such kernel,
with
\[
 \left|\sum_s b_s(u_x/s)\right|\geq c^k,\qquad
 \sum_s|b_s|^2\leq\frac{k!}{J_0^k},
\]
where $c>0$ is fixed and the support consists of products of $k$
distinct primes of $\mathcal P_0$.
Indeed each such product has coefficient
$k!J_0^{-k}\prod_{p\mid s}\epsilon'_p$, and there are
$\binom{J_0}{k}$ of them.
With $S=(2Z)^k$, we have $S=X^{1+o(1)}$ and
$S>mX+|h|$ for large $T$, while
\begin{equation}\label{eq:quad-coefficient-budget}
 S\frac{k!}{J_0^k}
 \leq\exp\bigl(O_\delta(\log Z+k\log T)\bigr)=X^{o(1)}.
\end{equation}
Apply the signed version of \eqref{eq:quad-heath-brown}, with
$U=mX+|h|$.  Since $c^{-2k}=X^{o(1)}$, the number of distinct
kernels on these subsets is at most
\begin{equation}\label{eq:quad-kernel-count}
 X^{3\varepsilon+o(1)}
\end{equation}
for every fixed $\varepsilon>0$; the estimate with exponent
$2\varepsilon+o(1)$ would also suffice.

\paragraph{Populations in a kernel and disjoint prime supports.}
For a fixed kernel $u$, the equality $ut_x^2=mx-h$ and
$(h,m)=1$ give $(ut_x,m)=1$.
The unit congruence $ut^2\equiv-h\pmod m$ has at most
$O(2^{\omega(m)})$ roots: there are at most two at an odd prime
power, at most four at a power of two, and the Chinese remainder
theorem applies.  On either of our endpoint intervals the range of
$t_x$ has diameter at most $\sqrt{mX/|u|}$, because the range of
$t_x^2$ has length at most $mX/|u|$ and
$|\sqrt a-\sqrt b|\leq\sqrt{|a-b|}$ for $a,b\geq0$.
Counting the possible roots in their residue classes gives
\begin{equation}\label{eq:src14}
 \begin{aligned}
 |\{x\text{ on a fixed side}:u_x=u\}|
 &\ll 2^{\omega(m)}\left(\frac{\sqrt{mX/|u|}}m+1\right)\\
 &\ll\sqrt{X/|u|}+X^{o(1)}.
 \end{aligned}
\end{equation}
Here $2^{\omega(m)}\ll\sqrt m$ uniformly: in the product
$\prod_{p\mid m}2/\sqrt p$, only the factors at $2$ and $3$ exceed
one.  Also $m=X^{o(1)}$.

Let $0<c_1\leq1/4$ be an absolute constant for the combined
zero-free region and Landau--Page statement recalled below, and put
$C_{\mathrm z}=1+c_1^{-1}\geq5$.
Fix, in this order,
\[
 0<\eta\leq\min\{1/1000,1/(200C_{\mathrm z})\},\qquad
 0<\varepsilon<\eta/100
\]
in
\eqref{eq:quad-kernel-count}.
The kernels with $|u|>X^\eta$ contribute at most
$X^{3\varepsilon+o(1)}(X^{1/2-\eta/2}+X^{o(1)})$
points on either side.  This is negligible compared with
$\sqrt X/(\log X)^3$.
We retain sets $\mathcal U\subset A_X$ and $\mathcal V\subset D_X$,
each of size $\gg_\delta\sqrt X/(\log X)^3$, on which all kernels
are nonzero and have absolute value at most $X^\eta$.

No prime divides a kernel on both sides.  In fact, if a prime divides
$v=u_x$ and $w=u_y$ for $x\in\mathcal U$, $y\in\mathcal V$, then
it divides $m(x-y)$ and is coprime to $m$, hence divides $x-y$.
But $x-y\in A+B$ is an eventual prime exceeding $X^{9/10}$,
whereas that divisor is at most $X^\eta$.
Taking $\eta<9/10$ makes this impossible.
Thus every such product $vw$ is signed squarefree.  Moreover its
prime factors determine $|v|$ and $|w|$: the prime supports of all
kernels on the first side and all kernels on the second side are
disjoint.  Only the bounded choice of signs remains.

\pdfbookmark[2]{Exceptional characters and the collision contradiction}{quadratic.collision}
\paragraph{Exceptional characters and split primes.}
We recall the precise unconditional input about Dirichlet
$L$-functions that is needed here.  For the absolute constant $c_1$
fixed above, among primitive nonprincipal characters of conductor at
most $Q$, all zeros with imaginary part of modulus at most one
satisfy
$\Re\rho\leq1-c_1/\log Q$, apart from at most one real simple
zero of a real character.  This is the classical zero-free region
together with the Landau--Page theorem.  For every fixed
$\varepsilon_1>0$, such a real zero satisfies
$1-\beta\gg_{\varepsilon_1}q^{-\varepsilon_1}$, by Siegel's
theorem, where $q$ is its conductor.  These statements are used only
for existence and asymptotics; no effective constant from Siegel's
theorem is required.  We use the zero-free region and Page theorem in
\cite[Theorem~11.3 and Corollary~11.10]{MontgomeryVaughan2007},
and the real-zero bound in
\cite[Corollary~11.15]{MontgomeryVaughan2007}.

Apply them with $Q_0=4X^{2\eta}$.
The primitive quadratic character corresponding to a signed
squarefree integer $vw\ne1$ has conductor at most $4|vw|\leq Q_0$.
If the exceptional character has conductor
$q>(\log X)^{100}$, discard the pairs $(x,y)$ that produce it.
The character determines $vw$ and hence the two kernel magnitudes
by the preceding disjointness.  Equation~\eqref{eq:src14} bounds
the number of discarded pairs by
\[
 O\left(\frac X{\sqrt{|vw|}}+X^{1/2+o(1)}\right)
 \ll\frac X{(\log X)^{50}}+X^{1/2+o(1)}.
\]
This is negligible compared with
$|\mathcal U||\mathcal V|\gg_\delta X/(\log X)^6$.
There are $O(X^{2\eta})$ possible ordered pairs of signed kernels.
Among the remaining pairs we can therefore fix kernels $v,w$ whose
population product is
$\gg_\delta X^{1-2\eta}/(\log X)^6$.
Each population is at most $O(\sqrt X)$ by
\eqref{eq:src14}; consequently each is at least
$X^{1/2-3\eta}$ for sufficiently large $T$.
Let $\mathcal R_1,\mathcal R_2$ be the resulting sets of positive
roots $t_x,t_y$.  The root maps are injective, their respective sizes
are these populations, and their diameters are at most
$X^{1/2+o(1)}$.

Set $Q_2=X^{1/2-5\eta}$.  We claim that our choice of $\eta$ gives
\begin{equation}\label{eq:src15}
 \sum_{\substack{p\leq Q_2\\p\nmid2mvw\\(vw/p)=1}}
                  \frac{\log p}{p}\geq .03\log X.
\end{equation}
If $vw=1$, this follows directly from Mertens' estimate after
removing the prime divisors of $2m$.
Otherwise let $\chi$ be the primitive quadratic character attached
to $vw$, of conductor $q\leq Q_0$, and put
\[
 s=1+\frac{10}{\log X},\qquad s_0=1+\frac1{\log Q_0}.
\]
For our fixed $\eta$ and sufficiently large $X$,
$1<s\leq s_0\leq2$.
The completed-function Hadamard formula gives, for real $1<u\leq2$,
\begin{equation}\label{eq:quad-log-derivative}
 \sum_\rho\Re\frac1{u-\rho}
       =\frac{L'}L(u,\chi)+\frac12\log q+O(1),
\end{equation}
where the zeros are the nontrivial zeros, counted with multiplicity.
For clarity, the completed logarithmic derivative is
$L'/L(u,\chi)+\tfrac12\log(q/\pi)
+\tfrac12\Gamma'/\Gamma((u+\kappa)/2)$, with
$\kappa\in\{0,1\}$; the real part of its Hadamard constant cancels
the sum of $\Re(1/\rho)$.  The gamma term is bounded on this
interval.  This proves precisely \eqref{eq:quad-log-derivative},
with a uniform $O(1)$; see also
\cite[Section~3.7.4, Equations (3.97)--(3.102)]{Helfgott}.

The zero sum is positive.  At $s_0$ it is at most $2\log Q_0$ for
sufficiently large $X$, because
the Euler series bounds $|L'/L(s_0,\chi)|$ by
$-\zeta'/\zeta(s_0)=\log Q_0+O(1)$ and $q\leq Q_0$.
For a nonexceptional zero $\rho=\beta+i\gamma'$ with
$|\gamma'|\leq1$, set $a=s-\beta$ and $b=s_0-\beta$.
The zero-free region gives $a\geq c_1/\log Q_0$, while
$0\leq b-a\leq1/\log Q_0$.  Hence
\[
 \frac{a/(a^2+(\gamma')^2)}{b/(b^2+(\gamma')^2)}
 \leq\frac ba\leq1+c_1^{-1}.
\]
If $|\gamma'|>1$, the same ratio is at most five, since
$0<a\leq b\leq2$.  Thus the sum of all nonexceptional terms at $s$
is at most $2C_{\mathrm z}\log Q_0$, with a constant independent of
$\eta$.
If the exceptional zero is a zero of the retained character, its
conductor is at most $(\log X)^{100}$.  Siegel's estimate with
$\varepsilon_1=1/200$ then gives
\[
 (s-\beta)^{-1}\leq(1-\beta)^{-1}
          \ll q^{1/200}\leq(\log X)^{1/2}=o(\log X).
\]
Using \eqref{eq:quad-log-derivative} and removing the absolutely
convergent prime-power terms yields
\begin{equation}\label{eq:quad-signed-prime-mass}
 \sum_p\frac{\chi(p)\log p}{p^s}
       =-\frac{L'}L(s,\chi)+O(1)
       \geq-2C_{\mathrm z}\log Q_0-o(\log X)
\end{equation}
with the fixed absolute constant $C_{\mathrm z}$.

On the other hand,
\[
 \sum_p\frac{\log p}{p^s}=\frac1{10}\log X+O(1),\qquad
 \sum_{p>Q_2}\frac{\log p}{p^s}
       =\frac1{10}e^{-5+50\eta}\log X+O(1).
\]
The second equality follows by partial summation from
$\sum_{p\leq y}\log p/p=\log y+O(1)$.
The primes dividing $2mvw$ have total $\log p/p$-weight
$O(\log\log X)$: primes up to $\log X$ satisfy this by Mertens,
and those above it contribute at most
$\log|2mvw|/\log X=O(1)$.
For other primes, $\ind_{\chi(p)=1}=(1+\chi(p))/2$.
Combining these facts with \eqref{eq:quad-signed-prime-mass}, the
split-prime mass with denominator $p^s$ is at least
\[
 \bigl(.05-2C_{\mathrm z}\eta-.1e^{-5+50\eta}-o(1)\bigr)\log X.
\]
Our preceding choice gives $2C_{\mathrm z}\eta\leq.01$ and
$.1e^{-5+50\eta}<.001$, so this coefficient exceeds $.03$ for
sufficiently large $X$.  Replacing $p^s$ by $p$ only increases the
split-prime sum and proves \eqref{eq:src15}.

\paragraph{The final collision contradiction.}
For $i=1,2$ let $\beta_i(p)$ be the probability that two independent
uniform elements of $\mathcal R_i$ are congruent modulo $p$.
For unequal roots, the sum of $\log p$ over prime divisors of their
difference is at most the logarithm of the diameter.
The equal-root contribution is
$O(Q_2/|\mathcal R_i|)=O(X^{-2\eta})$.
Consequently
\begin{equation}\label{eq:quad-final-collision-upper}
 \sum_{p\leq Q_2}\log p\bigl(\beta_1(p)+\beta_2(p)\bigr)
                       \leq(1+o(1))\log X.
\end{equation}
Always $\beta_i(p)\geq1/p$.
For a prime counted in \eqref{eq:src15}, choose a square root
$c_p$ of $w/v$ modulo $p$.  The occupied residue sets of
$\mathcal R_1$ and $c_p\mathcal R_2$ are disjoint.
Indeed, equality of residues would give
$vt_x^2\equiv wt_y^2\pmod p$, hence $p\mid m(x-y)$.
Since $p\nmid m$ and $x-y$ is a prime exceeding $X^{9/10}>Q_2$,
this is impossible.
If the two occupied sets have cardinalities $a_p,b_p$, their
collision probabilities are at least $1/a_p,1/b_p$, and
$a_p+b_p\leq p$.  Multiplication by $c_p$ preserves the second
collision probability, so
$\beta_1(p)+\beta_2(p)\geq4/p$ at every split prime under
consideration.
Thus the left side of \eqref{eq:quad-final-collision-upper} is at
least
\[
 2\sum_{p\leq Q_2}\frac{\log p}{p}
 +2\sum_{\substack{p\leq Q_2\\p\nmid2mvw\\(vw/p)=1}}
          \frac{\log p}{p}
 \geq(1-10\eta+.06+o(1))\log X.
\]
Our choice of $\eta$ contradicts
\eqref{eq:quad-final-collision-upper}.  This proves the proposition.
\end{proof}

\begin{corollary}\label{cor:quad-harmonic}
For fixed $0<\alpha<\beta$, as $L_*\longrightarrow\infty$,
\[
 \sum_{\alpha L_*\leq\log\log p\leq\beta L_*}\frac1p
       \max_{t\in\F_p}
       \left|\E_{x\in S_p}\chi_p(x-t)\right|=o(L_*).
\]
\end{corollary}
\begin{proof}
Partition the range of $\log p$ into dyadic intervals $[T,2T]$.
Equation~\eqref{eq:src3}, divided by $T$, bounds the contribution of
each interval with weight $1/p$ by a quantity tending to zero
uniformly for $T\geq\tfrac12e^{\alpha L_*}$.
There are $O(L_*)$ intervals, including at most two truncated ones;
positivity permits enlarging the latter to full intervals.  Their
sum is therefore $o(L_*)$.
\end{proof}
